\clearpage
\section{分析类模型}
\AddToHook{env/tabularx/begin}{\linespread{1.08}\selectfont}
\subsection{分析边界与责任划分}
本章从前文四个用例继续展开。课程材料第 280 页进入分析建模，第 287--288 页说明边界类、控制类和实体类，第 313--332 页讨论分析序列图及责任分配；第 398--403 页说明关系数据库映射。这里沿用这条顺序：先确定对象负责什么，再画对象间的调用，最后设计表结构。

分析类图中的矩形按名称、属性和操作分栏。实线表示关联，端点数字表示多重性；实心菱形位于整体端，表示组成对象归整体管理。属性前的“/”表示派生值。分析阶段不把主键、外键或数据库访问操作加到业务实体中。

\begin{table}[H]
\centering\caption{用例与边界、控制、实体的对应关系}
\small
\begin{tabularx}{\textwidth}{L{1.3cm}L{3.3cm}L{3.4cm}Y}
\toprule
用例 & 边界类 & 控制类 & 主要实体类\\\midrule
UC-01 & 前台接口、消息服务接口 & 登记预约控制 & 患者、预约、医生、诊室、预约变更\\
UC-02 & 收费接口、支付服务接口 & 录入收费控制 & 诊疗、收费单、收费明细、收费项目、支付、欠款\\
UC-03 & 医生接口、技工所接口 & 录入义齿订单控制 & 诊疗、义齿订单、产品、技工所、订单牙位、订单进展、提醒\\
UC-04 & 负责人接口、前台接口、消息服务接口 & 查询统计与提醒控制 & 预约、收费单、义齿订单、提醒\\
维护与诊疗 & 管理接口、医生接口 & 基础信息维护控制、录入诊疗控制 & 医院诊所、基础选项、员工岗位、价目、产品、诊疗\\
\bottomrule
\end{tabularx}
\end{table}

前台接口等人机边界接收输入并显示结果；支付服务接口、技工所接口和消息服务接口适配外部交互。控制类按用例协调对象，不包揽计价、欠款计算和订单阶段判断。后面序列图中的类名均沿用本表；名字前的冒号表示该类的一个对象。患者查询按患者档案及登记日期检索，基础信息维护和录入诊疗为前文功能清单中的配套用例，复用相同实体。

\clearpage
\subsection{患者、预约与诊疗}
患者负责核对身份，预约负责确认时段及改变预约阶段，诊疗保存医生实际完成的诊疗内容。医生和诊室是预约需要占用的资源，冲突检查要同时考虑两者。这里的“医生”是员工承担的业务角色，数据库阶段仍引用员工表，不另造一套人员档案。

\begin{figure}[H]\centering
\reportgraphic{\textwidth}{0.72\textheight}{07_classes_appointment.pdf}
\caption{患者、预约与诊疗分析类图}\label{fig:analysis-appointment}
\end{figure}

一位患者可以有多次预约和诊疗。一次预约最多对应一次诊疗，现场到诊可以没有预约。改约或取消保留原时段、原阶段、原因、操作人和发生时间。

预约时段包含开始时刻，不包含结束时刻，前一预约结束时可接下一预约。“新开始早于旧结束，且新结束晚于旧开始”表示重叠。已取消和爽约记录不占用未来时段。

\clearpage
\subsection{收费、支付与欠款}
收费单统管自己的收费明细和支付尝试。明细保存开单时采用的名称、单位与单价，后续调整价目不会改变已开单据。支付对象确认外部交易结果；收费单据此计算有效已收，欠款对象跟踪未结清款项。

\begin{figure}[H]\centering
\reportgraphic{\textwidth}{0.72\textheight}{08_classes_billing.pdf}
\caption{收费、支付与欠款分析类图}\label{fig:analysis-billing}
\end{figure}

应收等于明细金额之和减去整单合法折扣。已收只累计成功的收款，扣除成功退款；失败和待确认支付不改变已收。欠款等于应收减有效已收，不能用“收到回调”代替“支付成功”。

每张收费单可以多次尝试支付或分次收款。首次欠款时建立欠款对象，结清后保留结清时间；从未发生欠款的收费单可以没有欠款对象。

\clearpage
\subsection{义齿订单与交付协作}
义齿订单负责校验修复要求、接受合法回执和决定当前阶段，订单进展负责留存每次变化。一个订单只选用一种产品，可以包含多个牙位；不同产品分开建立订单，便于分别约定交期和返工。订单附件是订单的组成部分，文件名、格式和存储位置在数据库设计中展开。

\begin{figure}[H]\centering
\reportgraphic{\textwidth}{0.72\textheight}{09_classes_order.pdf}
\caption{义齿订单、进展与提醒分析类图}\label{fig:analysis-order}
\end{figure}

订单至少保留一条进展，第一条记录草稿建立。技工所回执必须指向已有订单，来源与事件号用于识别重复通知。试戴不合格时追加返工进展，保留原来的试戴结果；已交货订单不能直接改回“已接收”。

\clearpage
预约阶段、支付结果和订单阶段都有明确的业务含义，因此在模型中保留。它们的改变由实体操作完成，控制类不能任意修改这些属性。类图沿用课程的三栏矩形、可见性符号和操作签名；返回类型 Patient、Bill、OrderCopy 等分别表示患者、收费单和订单副本，Metric 表示带样本量的指标结果。图中未强加继承关系，员工的多个岗位以角色关联表达。

\begin{table}[H]\centering\caption{关键操作的责任归属}
\small
\begin{tabularx}{\textwidth}{L{3.0cm}L{4.4cm}Y}
\toprule
对象 & 对外操作 & 负责判断的业务规则\\\midrule
预约 & 确认、改约、取消、登记到诊、汇总履约 & 医生与诊室的时段占用、有效阶段及预约兑现口径\\
收费单 & 拟定、计算应收、应用支付结果、汇总回收 & 价格快照、合法折扣、有效已收、欠款和回收口径\\
支付 & 登记尝试、确认结果 & 请求是否重复、结果能否迁移、成功金额是否已计入\\
义齿订单 & 保存草稿、发送、受理回执、登记试戴、确认交货、申请返工 & 必填修复要求、阶段迁移、重复回执和交付口径\\
提醒 & 排定、判断到期、确认发送 & 计划时间、收件人及重复发送限制\\
\bottomrule\end{tabularx}
\end{table}

\subsection{功能覆盖核对}
\begin{table}[H]\centering\caption{项目功能与模型覆盖}
\small\begin{tabularx}{\textwidth}{L{2.6cm}L{6.1cm}Y}
\toprule 功能域 & 对象与持久化数据 & 查询或处理结果\\\midrule
基础信息 & 医院诊所、员工岗位、产品、价目、分类选项 & 颜色、盒色、单位、订单类型及附件配置可维护、可停用\\
患者与诊疗 & 患者、诊疗；登记时间、接诊医生和主诉 & 按身份、电话、登记期间查患者，回看诊疗历史\\
预约 & 预约、预约变更、医生与诊室关联 & 今日、明日、期间查询，冲突检查和改约留痕\\
费用 & 收费明细、支付、欠款 & 项目计价、收款与退款、欠款查询及单据追溯\\
义齿订单 & 产品、牙位、颜色、附件、进展、交期 & 试戴、已交货、返工和逾期查询\\
统计与提醒 & 业务记录汇总、预约或订单提醒 & 指标、样本、异常明细、通知计划及发送结果\\
\bottomrule\end{tabularx}\end{table}

\clearpage
\section{分析序列模型}
\subsection{UC-01 登记患者并预约诊疗}
前台作为外部执行者发起调用，其余参与对象是软件内部对象。执行者采用课件中的普通人形符号。软件对象的侧边竖线圆表示边界类，带开放箭头圆表示控制类，底线圆表示实体类。同步调用采用实线和直边实心三角箭头，返回采用虚线和开放箭头，细长矩形表示执行片段。消息箭头落在执行条边缘，人物、对象和连线采用统一的细黑线。

图中从前台提交要求开始，省略窗口打开等操作。预约对象在提交时形成，确认前不占用有效预约名额；身份未核对清楚时先返回待核对结果，不强行合并患者档案。

\begin{figure}[H]\centering
\reportgraphic{\textwidth}{0.70\textheight}{10_sequence_appointment.pdf}
\caption{UC-01 登记与预约分析序列图}\label{fig:seq-appointment}
\end{figure}

\clearpage
\subsection{UC-02 录入收费并处理欠款}
收费单、明细和待确认支付先在本地保存，随后才向支付服务发起请求。这使超时后仍能按原请求号核对结果。下图展开基本路径；守卫条件和图下说明补足成功、失败、待确认三种分支。

\begin{figure}[H]\centering
\reportgraphic{\textwidth}{0.80\textheight}{11_sequence_billing.pdf}
\caption{UC-02 收费、支付与欠款分析序列图}\label{fig:seq-billing}
\end{figure}

\clearpage
\subsection{UC-03 录入义齿订单并追踪交付}
\subsubsection{医生建立并发送订单}
诊疗对象先确认业务依据，义齿订单再核对产品、牙位、数量、颜色和必要附件。发送订单与技工所接受订单是两个结果，不能把一次发送成功当作技工所已经开始加工。

\begin{figure}[H]\centering
\reportgraphic{\textwidth}{0.80\textheight}{12_sequence_order_send.pdf}
\caption{UC-03 订单建立与发送分析序列图}\label{fig:seq-order-send}
\end{figure}

\clearpage
\subsubsection{技工所回执与进展追踪}
本图承接\figref{fig:seq-order-send} 的订单号，展开稍后到达的回执。技工所接口核对来源和报文格式，订单对象判断事件是否重复、阶段变化是否合法。只有处理成功后才确认回执，失败时保留重试或核对入口。

\begin{figure}[H]\centering
\reportgraphic{\textwidth}{0.80\textheight}{13_sequence_order_receipt.pdf}
\caption{UC-03 技工所回执分析序列图}\label{fig:seq-order-receipt}
\end{figure}

\clearpage
\subsection{UC-04 查询统计与订单提醒}
控制类向相应实体请求按统一期间汇总，不逐字段取出数据后自行判断业务含义。预约负责兑现口径，收费单负责欠款回收口径，义齿订单负责交期口径。返回结果都带分子、分母、样本量及异常明细，便于负责人核对。

\begin{figure}[H]\centering
\reportgraphic{\textwidth}{0.80\textheight}{14_sequence_statistics.pdf}
\caption{UC-04 统计与待处理订单分析序列图}\label{fig:seq-statistics}
\end{figure}

\clearpage
\section{数据库逻辑设计}
\subsection{从分析类到关系表}
课程在关系数据库部分区分类、属性和对象与表、列和行的对应关系，并把一对多关联的外键放到多的一端。本设计采用关系数据库逻辑模型，字段类型暂用通用 SQL 名称，具体数据库选型留到实现阶段。

数据库图参考 Shopify 和 Quickbooks 的 ERD：用表框逐行列出字段、键和外键目标，以业务分区组织主表、明细、支付与历史记录。下图先给出完整关系，再分别展开预约、费用和义齿业务。全库共 20 张表、44 个外键，覆盖基础资料、患者、诊疗、预约、收费、订单及通知；统计直接使用这些业务记录。

\begin{table}[H]\centering\caption{分析对象到关系表的映射}
\small\begin{tabularx}{\textwidth}{L{3.0cm}L{5.4cm}Y}
\toprule
分析对象 & 关系表 & 映射理由\\\midrule
患者、医生、诊室 & patient、employee、clinic\_room & 医生以员工角色参与，患者有独立档案\\
基础信息与员工岗位 & clinic、basic\_option、employee\_role & 诊所资料、分类选项与员工兼岗分别维护\\
预约、预约变更、诊疗 & appointment、appointment\_event、encounter & 当前安排与变更历史分开，诊疗可选关联预约\\
收费单、收费明细、收费项目 & bill、bill\_line、price\_item & 明细引用价目并保留成交快照\\
支付、欠款 & payment、bill；欠款查询视图 & 保存每次支付尝试，欠款按有效收款计算，催收日期归属收费单\\
产品、技工所、义齿订单 & product、dental\_lab、dental\_order & 订单引用一个产品和一家技工所\\
牙位、进展、附件、提醒 & order\_tooth、order\_event、order\_attachment、reminder & 多值结构拆开保存，避免逗号拼接牙位或覆盖历史\\
\bottomrule\end{tabularx}\end{table}

图中 PK 为主键，FK 为外键，UK 为唯一约束，问号表示允许空值。UK* 表示联合唯一约束，组合范围在后文列出。每个字段单列一行，并在右侧写明外键引用的表；连线箭头从外键所在表指向被引用表。同一对表的多个外键合并为一条连线，各字段的具体目标仍逐一保留，表\ref{tab:foreign-keys} 列出完整清单。

订单牙位与员工岗位使用复合主键，其余表使用不带业务含义的 id 作为主键，业务编号另设唯一约束。医生、诊室等对象关联转成外键；边界类和控制类不直接转成表。欠款是收费单的派生业务对象，数据库通过查询视图计算余额，不再单设欠款表。形成时间、到期日、结清时间和催收备注保存在收费单中。

\clearpage
\subsection{完整数据库设计图}
\figref{fig:erd-overview} 将 20 张表放在同一张图中。患者、预约、诊疗、收费单和义齿订单构成业务主链；员工、诊室、价目和产品提供基础资料，变更、支付尝试、回执、附件和提醒保留处理过程。颜色用于区分业务分区，字段名称、类型和关系均保留为矢量内容。

\begin{figure}[H]\centering
\reportgraphic{\textwidth}{0.68\textheight}{15_erd_overview.pdf}
\caption{口腔诊所完整数据库设计图（20 张表）}\label{fig:erd-overview}
\end{figure}

全图以独立矢量 PDF 嵌入报告，页面中按比例显示，放大后可以逐字段查看。后三张局部图用于阅读密集区域，重复出现的引用表仅列主键，不增加实际表数。

预约除了起止时刻，还保存创建、确认、到诊、完成及取消时刻，用于核对确认周期和兑现结果。收费单保存折扣原因及催收信息；支付保存请求、结果、核对时间和失败原因，支持分次收款、超时核对与退款。订单保留发送、接收、试戴和交货时刻，进展表保留事件来源与事件号，避免只凭当前阶段倒推业务过程。

基础选项集中维护岗位、单位、颜色、盒色、订单类型、预约类型、支付方式、通知渠道和附件类型。相同选项表可以被多个业务字段引用，但业务操作必须核对其类别。例如，\texttt{role\_id} 只能引用岗位类别，不能引用颜色或单位。

\clearpage
\subsection{患者、预约与诊疗 ERD}
患者保存登记时间和经办人，支持期间登记查询。预约保存主诉、开始与结束时间，变更表保存变更前后值。诊疗可以没有预约，但必须属于一个患者并由一名医生完成。手机号可以多人共用，不能强制唯一。

\begin{figure}[H]\centering
\reportgraphic{\textwidth}{0.74\textheight}{16_erd_appointments.pdf}
\caption{患者、预约与诊疗数据库设计图}\label{fig:erd-appointment}
\end{figure}

\clearpage
\subsection{收费、支付与欠款 ERD}
收费单通过诊疗关联患者；明细可以关联同次诊疗下的义齿订单，并保存开单时采用的项目名称、单位、数量和单价。支付表同时记录收款、退款和未成功的尝试，退款通过 \texttt{refund\_of\_id} 回到原成功收款。

\begin{figure}[H]\centering
\reportgraphic{\textwidth}{0.70\textheight}{17_erd_billing.pdf}
\caption{收费、支付与欠款数据库设计图}\label{fig:erd-billing}
\end{figure}

欠款查询先分别按收费单汇总明细和有效支付，再将两个汇总结果关联，防止明细与支付直接连接造成金额重复。应收为明细数量乘成交单价、扣除行折扣和整单折扣；有效已收为成功收款减成功退款；欠款余额为两者之差。失败、待确认或已作废单据不计入已收或有效欠款。

\texttt{debt\_opened\_at} 记录首次形成欠款的时刻，\texttt{due\_date} 记录约定到期日，\texttt{settled\_at} 记录最近一次结清时刻。退款使欠款重新产生时清除结清标记，并保留原支付记录。查询期初欠款时按截止时刻重算历史收付，不能只读取今天的结清标记。

\clearpage
\subsection{义齿订单与进展 ERD}
订单按牙位和进展拆分多值结构。产品名称、材料等主数据由产品表维护；订单上的颜色、数量、单位和承诺日期描述本次委托。草稿可尚无牙位，发送前必须至少有一个有效牙位；图中同时展开附件和提醒。

\begin{figure}[H]\centering
\reportgraphic{\textwidth}{0.70\textheight}{18_erd_orders.pdf}
\caption{义齿订单、牙位与进展数据库设计图}\label{fig:erd-order}
\end{figure}

订单附件独立保存文件位置、类型、大小、校验值和上传人；文件字节放在文件存储中。提醒记录目标、收件人、渠道、内容、计划时间、重试次数与发送结果。每条提醒恰好关联一个预约或一张义齿订单，不能同时关联两者，也不能没有目标。

\clearpage
\subsection{外键、唯一性与数据约束}
下表列出完整 ERD 中全部外键。每一行左侧的多个字段分别引用右侧相应主键；同表内出现多种员工角色时，仍使用不同字段保存经办人与医生。字段、图和本表由同一份逻辑模型生成。

\begingroup\linespread{1}\fontsize{10.5}{13}\selectfont\renewcommand{\arraystretch}{1.05}
\begin{longtable}{L{8.0cm}L{6.5cm}}
\caption{外键引用清单}\label{tab:foreign-keys}\\
\toprule 外键字段 & 引用目标\\\midrule\endfirsthead
\toprule 外键字段 & 引用目标\\\midrule\endhead
clinic\_room.clinic\_id & clinic.id\\
employee.clinic\_id & clinic.id\\
employee\_role.employee\_id & employee.id\\
employee\_role.role\_id & basic\_option.id\\
patient.created\_by & employee.id\\
appointment.patient\_id & patient.id\\
appointment.doctor\_id、created\_by & employee.id\\
appointment.room\_id & clinic\_room.id\\
appointment.type\_id & basic\_option.id\\
appointment\_event.appointment\_id & appointment.id\\
appointment\_event.changed\_by & employee.id\\
encounter.patient\_id & patient.id\\
encounter.clinic\_id & clinic.id\\
encounter.doctor\_id & employee.id\\
encounter.appointment\_id & appointment.id\\
price\_item.unit\_id & basic\_option.id\\
bill.encounter\_id & encounter.id\\
bill.issued\_by & employee.id\\
bill\_line.bill\_id & bill.id\\
bill\_line.price\_item\_id & price\_item.id\\
bill\_line.dental\_order\_id & dental\_order.id\\
payment.bill\_id & bill.id\\
payment.method\_id & basic\_option.id\\
payment.refund\_of\_id & payment.id\\
product.unit\_id & basic\_option.id\\
dental\_order.encounter\_id & encounter.id\\
dental\_order.product\_id & product.id\\
dental\_order.lab\_id & dental\_lab.id\\
dental\_order.created\_by & employee.id\\
dental\_order.color\_id、box\_color\_id、type\_id、unit\_id & basic\_option.id\\
order\_tooth.order\_id & dental\_order.id\\
order\_event.order\_id & dental\_order.id\\
order\_event.changed\_by & employee.id\\
order\_attachment.order\_id & dental\_order.id\\
order\_attachment.type\_id & basic\_option.id\\
order\_attachment.uploaded\_by & employee.id\\
reminder.order\_id & dental\_order.id\\
reminder.appointment\_id & appointment.id\\
reminder.channel\_id & basic\_option.id\\
reminder.created\_by & employee.id\\
\bottomrule\end{longtable}
\endgroup


\clearpage
\begin{table}[H]\centering\caption{唯一性、检查及索引设计}
\small\begin{tabularx}{\textwidth}{L{3.0cm}Y}
\toprule 对象 & 约束与使用目的\\\midrule
预约时段 & starts\_at $<$ ends\_at；索引 (doctor\_id, starts\_at) 和 (room\_id, starts\_at)。普通唯一索引不能排除部分重叠，需配合下文事务检查。\\
收费明细 & 联合唯一 (bill\_id, line\_no)；quantity $>$ 0，unit\_price $\geq$ 0；项目折扣不大于该行原价；整单折扣不大于行金额合计。\\
支付 & request\_no 唯一；external\_txn\_no 在非空时唯一；amount $>$ 0；kind 区分收款和退款；退款必须引用原成功收款，同单累计退款不超过原成功金额。\\
订单牙位 & 联合主键 (order\_id, tooth\_code)，同一订单同一牙位不能重复；已发送订单至少有一个有效牙位。\\
订单进展 & 联合唯一 (order\_id, source, event\_no)，隔离不同来源的事件号；索引 (order\_id, occurred\_at)。\\
基础选项与岗位 & 基础选项联合唯一 (kind, code)，防止同类别编码重复；员工岗位的 role\_id 引用岗位类别，联合主键 (employee\_id, role\_id)，允许员工兼岗。\\
提醒 & 对非空 order\_id 或 appointment\_id，分别与 kind、scheduled\_at、recipient 构成联合唯一；两个目标恰好有一个非空。索引 (delivery\_result, scheduled\_at)。\\
附件 & byte\_size $>$ 0；mime\_type 和大小按允许范围检查；文件通过 storage\_key 定位。\\
\bottomrule\end{tabularx}\end{table}

姓名、联系方式和诊断文本不参与重复请求识别。业务编号、支付请求号和外部事件号各自解决不同问题，不能用一次时间戳兼任这些标识。订单附件不把文件字节塞进订单字段，牙位也不保存成逗号分隔字符串。

\subsection{事务边界与异常处理}
\textbf{预约确认。} 先按固定顺序锁定涉及的医生和诊室资源，在同一事务中检查有效预约的时段重叠，再保存预约及变更。改约同时处理旧、新资源。两个前台并发提交时，后进入检查的一方必须看到前一方的提交结果；仅在界面上检查一次不能保证这一点。

\textbf{收费与回执。} 收费单和明细先原子保存，支付请求使用唯一请求号。应用回执时锁定收费单及相应支付，核对金额、外部交易号和原结果，确认这次成功是否已计入，再更新收费单的欠款形成或结清标记。外部请求放在数据库事务之外；超时保留待确认，按同一请求号查询或重试。

\textbf{订单进展。} 同一事务内校验原阶段、保存新进展和改变订单当前阶段。数据库中的 stage 是供查询使用的当前值，必须与最新进展一致。技工所回执重复到达时返回已处理结果；待发送订单和待发送提醒由已保存记录驱动重试，不依赖窗口一直打开。

\textbf{跨表业务一致性。} 诊疗引用预约时，患者和医生必须与该预约一致；收费明细引用义齿订单时，该订单必须属于收费单所对应的诊疗。外键只能保证目标存在，这两条约束还需由业务操作在事务中核对。医生外键所指员工必须具备医生角色，停用的产品、技工所和诊室不能用于新业务。

\textbf{历史保留。} 已开单、已发送订单和成功支付保留原记录，通过退款、取消或返工追加原因。被业务表引用的患者、员工、价目和产品不直接删除，用停用标记限制后续选择。显示当前结果时可读取最新阶段，追溯时按进展或变更时间还原。

\subsection{分析与设计的追踪核对}
\begin{table}[H]\centering\caption{用例、对象协作与数据库设计追踪}
\small\begin{tabularx}{\textwidth}{L{1.3cm}L{4.2cm}L{4.1cm}Y}
\toprule 用例 & 分析序列图 & 核心表 & 核对场景\\\midrule
UC-01 & \figref{fig:seq-appointment} & patient、appointment、appointment\_event、encounter & 并发冲突、身份待核对、改约留痕、无预约到诊\\
UC-02 & \figref{fig:seq-billing} & bill、bill\_line、payment & 分次收款、待确认、重复回执、退款后余额\\
UC-03 & \figref{fig:seq-order-send}、\figref{fig:seq-order-receipt} & dental\_order、order\_tooth、order\_event、order\_attachment、reminder & 缺少牙位、发送重试、重复回执、试戴返工\\
UC-04 & \figref{fig:seq-statistics} & appointment、bill、payment、dental\_order、reminder & 零样本、口径追溯、逾期明细、提醒结果\\
\bottomrule\end{tabularx}\end{table}

\begin{table}[H]\centering\caption{统计口径与已有记录的对应}
\small\begin{tabularx}{\textwidth}{L{3.3cm}Y}
\toprule 指标 & 计算口径\\\midrule
预约兑现率 & 以预约开始时间落在统计期间的已确认、已到诊、已完成和爽约预约为分母，以其中已到诊或已完成数为分子；取消和待确认不纳入。\\
欠款回收率 & 固定期初欠款收费单队列，分母为期初未收余额，分子为该队列在期间实际回收的金额，扣除相关成功退款；期间新增欠款另列。\\
订单按期交付率 & 分母为承诺交期落在已结束统计期间的已发送且未取消订单，分子为其中实际交货不晚于承诺日期的订单；到期未交付仍计入分母。\\
记录关联完整度 & 分母为期间纳入核查的收费和订单记录，分子为能关联有效患者、诊疗及必要明细的记录；异常单独列出。\\
\bottomrule\end{tabularx}\end{table}

比率分母为零时显示“无可计算样本”。实现前还需由诊所确认允许的折扣与退款权限、哪些修复类型必须附带文件，以及提醒提前多久、通过什么渠道发送。这里将其保留为配置条件，未填入未经确认的费率或日期。
